## TEMP holding note — RQ2b attribution: NLME + Elastic Net + XGBoost on mean_cr_residual (2026-08-11)

**Cluster fit ran**: 2026-08-10 (`step1_rq2_fit.sh`, 15 jobs — 3 sets x 5 folds, 4survey only).
**Note**: Set3/Set4 initially FAILED on the cluster (`ElasticNetCV` crashed on NaN input — Set2's
10 columns happened to have zero missingness, so this wasn't caught by pre-cluster local
smoke-testing; Set3/4 pull in columns with a small nonzero missing fraction, e.g. `soilgrids_ph`,
the CEH layers). Fixed same day (`drop_rows_with_missing_features()` added to both
`run_rq2_attribution.py` and `evaluate_rq2_attribution.py`, applied identically to train and test
so NLME/EN/XGBoost all see the same rows) and the 10 affected fold-jobs resubmitted and completed
clean. **Synced + evaluated locally**: 2026-08-11 (`step1_rq2_sync.sh` then
`step1_rq2_evaluate.sh`). Not yet added to the main results artifact.

### How these numbers were generated

- **Fit** (cluster, CPU): NLME (`fit_nlme`, diagnostics saved directly, no held-out step —
  see `run_rq2_attribution.py`'s own header for why) + Elastic Net + XGBoost, on
  `nested_set2_top10`/`nested_set3_gated_terrain_wind_vif`/`nested_set4_gated_all_vif`
  (VIF-screened per `documentation/methodlogy_env_setpick.md`), `spatial_block_kfold`, split
  seed 42, 4survey only (6survey's 47 compartments too few for this — see
  `documentation/experiment_log.md`'s cohort-justification entry).
- **Evaluate** (local): reloads the EN (`joblib`)/XGBoost (native `save_model`/`load_model`,
  not joblib — see the 2026-08-10 cross-version pickling fix) checkpoints, re-derives the
  identical test split, scores both. NLME has no held-out evaluate step — its numbers below are
  computed directly on its own training data at fit time.
- **Pooled R2** here (unlike RQ1) is computed directly from this note's own aggregation of all 5
  folds' `predictions.csv` (concatenated, one R2 over the whole population), since RQ2's evaluate
  script itself only prints per-fold numbers — not from a saved `kfold_summary.json` (RQ2 doesn't
  currently call `models.common.kfold_summary`).
- **Coefficient tables**: mean ± SD computed across the 5 already-saved per-fold
  `nlme_fixed_effects.json` / `elastic_net_coefficients.csv` files — this is the fold-stability
  check flagged as a "medium tier" plot idea earlier, done here as a table instead.

### Results — pooled + per-fold R2 (4survey only)

| Set | n rows (pooled) | EN pooled R2 | EN fold R2 (mean±SD) | XGB pooled R2 | XGB fold R2 (mean±SD) | NLME variance explained (mean±SD) |
|---|---:|---:|---|---:|---|---|
| Set2 (`nested_set2_top10`) | 71,766 | 0.359 | 0.352±0.030 | 0.326 | 0.318±0.034 | 0.204±0.066 |
| Set3 (`nested_set3_gated_terrain_wind_vif`) | 71,727 | 0.325 | 0.318±0.048 | 0.299 | 0.288±0.056 | 0.049±0.048 |
| Set4 (`nested_set4_gated_all_vif`) | 71,330 | **0.358** | 0.350±0.040 | **0.338** | 0.329±0.041 | 0.175±0.064 |

Note the row-count drop from Set2→Set4 (71,766→71,330, ~0.6%) is entirely the missing-value drop
from the fix above, not a data problem.

### NLME fixed-effect coefficients, mean ± SD across 5 folds (sorted by |mean|)

**Set2**: `CanopyCover` +1.839±0.080, `dist_to_road` -1.312±0.110, `time_since_thinning`
-1.254±0.501, `time_since_thinning_missing` -1.079±0.624, `slope_degrees` +0.999±0.202,
`chelsa_bio12_precip_mm` -0.785±0.669, `gwa_weibull_a_50m` -0.537±0.136, `gwa_weibull_k_50m`
-0.451±0.231, `eastness` +0.370±0.016, `recent_thinning_5yr` -0.328±0.249, `tas_mean`
+0.228±0.315, `chelsa_gdd5_degc` +0.188±0.276, `gwa_weibull_a_10m` -0.181±0.059,
`local_relief_500m` -0.174±0.178.

**Set3**: `CanopyCover` +1.854±0.073, `time_since_thinning` -1.487±0.486,
`time_since_thinning_missing` -1.397±0.602, `topex` +1.316±0.171, `slope_degrees` +0.894±0.226,
`local_relief_500m` -0.620±0.308, `whcl` -0.476±0.360, `recent_thinning_5yr` -0.398±0.243,
`gwa_weibull_k_50m` -0.362±0.225, `eastness` +0.342±0.047, `gwa_wind_speed_10m` -0.126±0.065,
`gwa_weibull_a_10m` -0.125±0.046, `ceh_twi` +0.079±0.068, `gwa_weibull_a_50m` -0.040±0.104.

**Set4**: `CanopyCover` +1.863±0.071, `time_since_thinning` -1.476±0.531, `chelsa_bio12_precip_mm`
-1.416±0.552, `time_since_thinning_missing` -1.323±0.622, `topex` +1.323±0.138, `dist_to_road`
-1.216±0.116, `slope_degrees` +0.861±0.213, `whcl` -0.459±0.358, `recent_thinning_5yr`
-0.404±0.255, `gwa_weibull_k_50m` -0.360±0.223, `gwa_weibull_a_50m` +0.342±0.059, `eastness`
+0.326±0.018, `dist_to_scpt_boundary` -0.145±0.064, `local_relief_500m` -0.123±0.228,
`gwa_weibull_a_10m` -0.110±0.050, `gwa_wind_speed_10m` -0.108±0.061, `soilgrids_ph`
+0.069±0.200, `ceh_twi` +0.054±0.042, `tas_mean` +0.015±0.321.

### Elastic Net coefficients, mean ± SD across 5 folds (sorted by |mean|)

**Set2**: `time_since_thinning_missing` -2.751±0.223, `time_since_thinning` -1.982±0.236,
`CanopyCover` +1.926±0.034, `dist_to_road` -1.374±0.170, `slope_degrees` +1.070±0.113,
`chelsa_bio12_precip_mm` -0.915±0.385, `recent_thinning_5yr` -0.866±0.225, `gwa_weibull_a_50m`
-0.719±0.252, `eastness` +0.516±0.059, `local_relief_500m` +0.348±0.068, `gwa_weibull_k_50m`
-0.175±0.141, `tas_mean` -0.151±0.371, `gwa_weibull_a_10m` -0.097±0.086, `chelsa_gdd5_degc`
+0.003±0.296.

**Set3**: `time_since_thinning_missing` -3.308±0.422, `time_since_thinning` -2.233±0.238,
`CanopyCover` +1.869±0.023, `gwa_weibull_a_50m` -1.360±0.421, `slope_degrees` +1.230±0.243,
`recent_thinning_5yr` -0.963±0.355, `whcl` -0.720±0.264, `eastness` +0.507±0.106,
`gwa_weibull_k_50m` -0.278±0.145, `local_relief_500m` -0.258±0.186, `gwa_wind_speed_10m`
-0.204±0.090, `topex` -0.175±0.166, `gwa_weibull_a_10m` -0.060±0.065, `ceh_twi` -0.049±0.174.

**Set4**: `time_since_thinning_missing` -2.999±0.164, `time_since_thinning` -2.141±0.194,
`CanopyCover` +1.943±0.034, `dist_to_road` -1.332±0.153, `slope_degrees` +0.993±0.194,
`recent_thinning_5yr` -0.977±0.245, `chelsa_bio12_precip_mm` -0.620±0.452, `whcl` -0.594±0.204,
`gwa_weibull_a_50m` -0.579±0.545, `eastness` +0.472±0.077, `tas_mean` -0.334±0.217,
`dist_to_scpt_boundary` -0.299±0.093, `local_relief_500m` +0.164±0.098, `gwa_weibull_k_50m`
-0.136±0.175, `soilgrids_ph` +0.123±0.166, `gwa_wind_speed_10m` -0.121±0.067, `gwa_weibull_a_10m`
-0.082±0.084, `topex` +0.063±0.265, `ceh_twi` -0.038±0.156.

### Headline pattern

**`CanopyCover` and the two thinning-history columns (baseline, not "environmental") dominate
every set, both models, with tight cross-fold SD relative to their mean** (e.g. EN's
`time_since_thinning_missing`: -2.75 to -3.31 depending on set, SD always <0.5) — the single
clearest, most stable signal in this whole table. This echoes the same pattern RQ3's category
attribution found (`TEMP_rq3_category_attribution_results_2026-08-10.tex`, now a pre-VIF
snapshot but the qualitative finding likely still holds): stand-structure variables carry
disproportionate weight relative to the purely environmental candidates. **Not yet checked**:
whether this is a genuine forestry mechanism or a confound via how `CanopyCover` is derived
(flagged, not investigated, earlier this session).

**`topex`/`windward_topex`-family wind exposure and `slope_degrees` are the most stable
environmental (non-baseline) signals** — both NLME and EN agree on direction and magnitude across
Set3/Set4, SD small relative to mean. **Several variables flip sign or have SD comparable to
their mean** (`chelsa_gdd5_degc`, `tas_mean`, `topex` in EN Set4, `soilgrids_ph`) — these should
be read as "not reliably different from zero across folds," not as settled findings.

**Set3 is the weakest set here** (lowest R2 for both models, lowest NLME variance explained) —
consistent with the earlier finding that Set3's category composition is terrain-heavy with thin
wind representation and zero soil/climate coverage.

### XGBoost SHAP — added 2026-08-12, all 3 sets, 4survey

New: `models/spatial_attribution/compute_rq2_shap.py` — cheaper than RQ3's SHAP script since
RQ2's fit step already saves `xgboost_model.json` per fold; this is a pure reload + compute, no
refitting at all. Reuses the same generic `compute_shap_values_for_columns()` built for RQ3, no
new SHAP code. Ran for all 3 sets x 5 folds x 4survey = 15 combos, all succeeded.

**Mean |SHAP| across all 5 folds, top 6 per set**:

| Set2 | Set3 | Set4 |
|---|---|---|
| CanopyCover 1.474 | CanopyCover 1.448 | CanopyCover 1.481 |
| dist_to_road 1.220 | whcl 0.861 | dist_to_road 1.099 |
| chelsa_bio12_precip_mm 0.955 | slope_degrees 0.817 | chelsa_bio12_precip_mm 0.932 |
| slope_degrees 0.681 | gwa_weibull_a_50m 0.618 | slope_degrees 0.733 |
| chelsa_gdd5_degc 0.547 | gwa_weibull_k_50m 0.613 | whcl 0.452 |
| tas_mean 0.416 | local_relief_500m 0.593 | topex 0.431 |

**`CanopyCover` is #1 by mean |SHAP| in every set** — matches the NLME/EN coefficient ranking
above (a real correctness check, not assumed) and confirms this is the single most consistent
finding across every attribution method tried on this target so far, not an artifact of any one
model family.

**Not yet done**: cross-method agreement heatmap (NLME vs. EN vs. XGBoost SHAP — all three now
exist, table not built), 5-seed reseed (no seed sweep has been run for RQ2 at all — this note
only ever reflects split seed 42), SHAP beeswarm plot (data ready, not plotted).

### CanopyCover reverse-causation check — baseline-only (Set1) comparison, added 2026-08-13

**Context**: the project's own data cheat-sheet (`documentation/data_exploration_gpkg/Forest LiDAR
& Inventory – Key Terms Cheat Sheet.md`, §3.2/§3.4/§13) flags `CanopyCover` as a "downstream
consequence of the growth process," not a genuine driver — confirmed it and `elev_percentile_95th`
come from the SAME LiDAR CSV/ALS flight (not independent surveys), and correlations with all 3
targets are moderate (Pearson 0.35-0.47), ruling out literal duplication (`GapFraction`, an exact
1-CanopyCover duplicate, correlates far higher) but not resolving the causal-direction question.
This is NOT the same category as the `Age` circularity bug (Age is algebraically inside the CR
formula; CanopyCover is not in any of the 3 target formulas) — a judgment call, not a proven leak.

**Test run**: `nested_set1_baseline` (the 4 stand-structure columns alone — `CanopyCover`,
`time_since_thinning`, `time_since_thinning_missing`, `recent_thinning_5yr`) had never actually
been fit for RQ2 before (only Set2/3/4, which already include baseline + environmental candidates
together). Widened `--set-name` choices in `run_rq2_attribution.py`/`evaluate_rq2_attribution.py`
(additive, no behavior change for existing sets) and ran all 5 folds locally (CPU-only, ~10s/fold).

**Result — see the comparison table above this section, now with a Set1 row.** Two findings:
1. Baseline alone reaches only R2=0.185-0.220 (EN/XGBoost) — environmental variables add a real,
   non-trivial +0.10 to +0.15 R2 on top (48-83% relative increase) — baseline is nowhere close to
   "the whole model."
2. NLME's variance-explained metric (specifically BETWEEN-COMPARTMENT/spatial variance, not
   overall row-level variance) is ~0 for baseline alone (0.016+/-0.078, one fold negative) vs.
   0.05-0.20 for the full sets — baseline has real row-level predictive power but explains almost
   none of the spatial pattern, while environmental variables explain a real share of it.

**Conclusion**: supports treating `CanopyCover`/thinning explicitly as baseline stand-structure
CONTROLS in the write-up (expected to dominate row-level R2, not itself the interesting
attribution finding), with the environmental variables' contribution beyond that baseline as the
actual headline — a defensible, evidence-backed reframe rather than "CanopyCover just happens to
win." No code/model change made (baseline stays in every set, as before) — this was a
write-up-framing decision, backed by a real ablation number instead of an assumption either way.
